% =====================================================================
%  EQUIPO 1 (GGHL) — T7 Computación paralela — 26/11/2026
%  Integrantes observados: Diego, Roger, Lombeida.
%  Las diapositivas citadas (1/5 a 5/5) son las de la parte de Diego.
%  PENDIENTE: poner los niveles (?) y los mismos números en \notasugerida.
% =====================================================================

\section{Grupo 1 (GGHL) --- T7 Computación paralela}

\identificacion
{Grupo 1 (GGHL): Diego, Roger y Lombeida}
{T7 --- Computación paralela (Apache Spark)}
{26/11/2026}
{08:40}
{09:10}

\begin{dimensiones}
	\dimension{D1 Cumplimiento de la guía (10)}
	{3}
	{Conceptos clave: presentados por Diego (diap.~1--2). Demostración en vivo:
		Diego ejecutó la práctica RTV con Spark \texttt{local[*]} (diap.~3);
		Lombeida presentó una práctica pequeña; Roger no presentó práctica.
		Comparación: solo entre 2 y 4 particiones de la misma práctica (diap.~3).
		Recomendación: no se observó una recomendación explícita.}
	
	\dimension{D2 Fundamentación conceptual (20)}
	{3}
	{Diego explicó: definición de computación paralela y su objetivo
		(diap.~1); diferencia entre ejecución secuencial, paralela y distribuida con
		el ejemplo de vehículos (diap.~1); arquitectura de Spark con driver, gestor
		del clúster, nodos y ejecutores (diap.~2); tareas en hilos y núcleos
		(diap.~2); particiones y las operaciones \texttt{parallelize()},
		\texttt{map()} y \texttt{collect()} (diap.~3). No se detectaron errores
		graves de la guía. La explicación de Roger quedó incompleta.}
	
	\dimension{D3 Desarrollo técnico (15)}
	{2}
	{Diego explicó el código y lo ejecutó: Spark en modo \texttt{local[*]}
		en una sola computadora (diap.~3), es decir, no en procesos o equipos
		distintos, y sin caso de fallo. Resultado: 2 particiones, 2/2 tareas,
		1,8195\,s; 4 particiones, 4/4 tareas, 3,4206\,s (dos mediciones por
		configuración). Se comparó una sola variable (número de particiones) con
		un criterio medible (tiempo), no dos alternativas tecnológicas. Roger no
		hizo demostración de código; la práctica de Lombeida fue pequeña y su
		explicación del código, escasa.}
	
	\dimension{D4 Análisis y recomendación (20)}
	{3}
	{Se reconoció una limitación: en la prueba, aumentar las particiones
		elevó el tiempo, y en un clúster también influyen el reparto y la
		ubicación de los datos (diap.~3). No se indicó cuándo usar y cuándo no
		usar la computación paralela ni se ligaron razones a un caso concreto.}
	
	\dimension{D5 Comunicación (25)}
	{3}
	{Duración: 30\,min (08:40--09:10), sin tiempo asignado para comparar con
		$T \pm 2$\,min. Expusieron los tres integrantes. Diego explicó su parte con
		claridad. A Lombeida no se le entendió bien por su entonación. No hubo
		preguntas, así que no se observó respuesta a una pregunta técnica.}
	
	\dimension{D6 Fuentes (10)}
	{2}
	{En las diapositivas de Diego, cada dato tiene su cita visible (diap.~1--3),
		incluida la tabla de tiempos (``registro de la práctica''). Las referencias
		están en formato IEEE con DOI (diap.~4--5): fuentes indexadas (Springer,
		\emph{IEEE Transactions}, \emph{The Computer Journal}) y documentación
     oficial de Apache Spark. El dato del 80\,\% de Fortune 500 (diap.~1)
     proviene de la página del propio proyecto. No se verificó en las diapositivas de Roger y Lombeida.}
\end{dimensiones}

% Pon aquí los mismos niveles de la tabla, en orden D1..D6 (solo números):
\notasugerida{3}{3}{2}{3}{3}{2}

\retroalimentacion
  {Por parte de Diego buen dominio del tema y de explicación}
  {Por parte de los compañeros Roger y Lombeida se recomienda profundizar en la explicación práctica.}

\clearpage
